\section{Models}\label{sec:models}
\begin{figure*}[t]
\centering
\begin{minipage}[b]{0.50\textwidth}\centering
\resizebox{\linewidth}{!}{\begin{tikzpicture}[x=1cm,y=1cm]
\tikzset{
  nb/.style={rectangle,rounded corners=3pt,draw=black!60,fill=blue!8,text width=7.2cm,align=center,minimum height=0.62cm,font=\scriptsize},
  hd/.style={rectangle,rounded corners=3pt,draw=black!60,fill=green!10,text width=2.15cm,align=center,minimum height=0.95cm,font=\scriptsize},
  ar/.style={-{Stealth[length=1.8mm]},thick,draw=black!65}
}
\node[nb] (in) at (0,0) {15 joint angles (z\,=\,0; hidden joints masked)};
\node[nb] (l1) at (0,-0.95) {Linear 15$\rightarrow$256, BatchNorm, GELU};
\node[nb,fill=blue!16] (rb) at (0,-1.9) {2 $\times$ ResBlock(256): (Linear--BN--GELU--Dropout 0.3) $\times 2$, identity skip};
\node[hd] (h1) at (-2.55,-3.35) {\textbf{Pose head}\\ 256$\rightarrow$128$\rightarrow P$};
\node[hd] (h2) at (0,-3.35) {\textbf{Form head}\\ 256$\rightarrow$64$\rightarrow 1$};
\node[hd] (h3) at (2.55,-3.35) {\textbf{Deviation head}\\ 256$\rightarrow$128$\rightarrow 15$};
\draw[ar] (in)--(l1); \draw[ar] (l1)--(rb);
\draw[ar] (rb.south) -- ++(0,-0.28) -| (h1.north);
\draw[ar] (rb.south) -- (h2.north |- rb.south) -- (h2.north);
\draw[ar] (rb.south) -- ++(0,-0.28) -| (h3.north);
\node[font=\scriptsize,text=black!70] at (-2.55,-4.25) {pose name};
\node[font=\scriptsize,text=black!70] at (0,-4.25) {form score $\in[0,1]$};
\node[font=\scriptsize,text=black!70] at (2.55,-4.25) {per-joint deviation};
% cascade
\node[nb,fill=orange!12,text width=7.2cm] (cas) at (0,-5.25) {\textbf{Cascade at serving time:} model $v4$'s pose head \textbf{names}; a second, newer model's form/``not mine'' output \textbf{vetoes} poses outside the vocabulary and scores form; angle bands give the joint-level cues};
\end{tikzpicture}
}\\[2pt]
{\footnotesize (a) The 3-head residual MLP and the cascade.}
\end{minipage}\hfill
\begin{minipage}[b]{0.47\textwidth}\centering
\resizebox{\linewidth}{!}{\begin{tikzpicture}[x=1cm,y=1cm]
\tikzset{
  nb/.style={rectangle,rounded corners=3pt,draw=black!60,fill=blue!8,text width=5.0cm,align=center,minimum height=0.62cm,font=\scriptsize},
  ar/.style={-{Stealth[length=1.8mm]},thick,draw=black!65},
  an/.style={font=\scriptsize,text=black!70,text width=2.7cm,align=left}
}
\node[nb] (in) at (0,0) {Window: 60 frames $\times$ 33 joints $\times$ $(x,y,z)$};
\node[nb,fill=green!10] (b1) at (0,-1.0) {ST-GCN block 1: graph conv $3\!\rightarrow\!64$ + temporal conv (9 taps)};
\node[nb,fill=green!10] (b2) at (0,-2.0) {Block 2: $64\!\rightarrow\!128$};
\node[nb,fill=green!10] (b3) at (0,-3.0) {Block 3: $128\!\rightarrow\!256$};
\node[nb] (gp) at (0,-4.0) {Global average over joints and time};
\node[nb] (fc) at (0,-5.0) {Linear 256$\rightarrow$128, LayerNorm, GELU, Linear 128$\rightarrow$9};
\node[nb,fill=orange!12] (out) at (0,-6.0) {8 poses + \textit{transition / unknown}};
\foreach \a/\b in {in/b1,b1/b2,b2/b3,b3/gp,gp/fc,fc/out}{\draw[ar] (\a)--(\b);}
\node[an] at (4.65,-0.0) {pelvis-centred, scaled by hip width, 25\,fps};
\node[an] at (4.65,-2.0) {each block: $\mathrm{GELU}(\mathbf{A}_{\mathrm{norm}}\,W\mathbf{X})$, residual skip};
\node[an] at (4.65,-5.0) {0.89\,M parameters};
\end{tikzpicture}
}\\[2pt]
{\footnotesize (b) The ST-GCN sequence model (``flow check'').}
\end{minipage}
\caption{The two learned models. $P$ is the number of pose classes. For $P\approx22$ the MLP has about 0.36\,M parameters; the nine-class ST-GCN has about 0.89\,M (both counted from the model definitions).}
\label{fig:models}
\end{figure*}

\subsection{3-head residual MLP}
The pose model takes the 15-angle vector (Section~\ref{sec:features}) and passes it through a 256-unit input layer (Linear, batch normalisation \cite{ioffe2015bn}, GELU \cite{hendrycks2016gelu}) and two residual blocks \cite{he2016resnet} of width 256 (two Linear--BN--GELU--Dropout stages with an identity skip). Three heads share this trunk (Fig.~\ref{fig:models}a): a \emph{pose head} (256$\to$128$\to P$ logits), a \emph{form head} (256$\to$64$\to1$, a correctness probability), and a \emph{deviation head} (256$\to$128$\to$15 per-joint deviations). Sharing the trunk makes one forward pass of about 0.36\,M parameters answer ``which pose, how well, which joints'' in a few milliseconds on a CPU.

\subsection{The cascade: one model names, another vetoes}\label{sec:cascade}
Through the original decision path, a rule engine overrode the MLP on every disagreement, so replacing the MLP with a better one changed nothing (identical to rules-only, Table~\ref{tab:policy}). We therefore split the roles. The \emph{namer} is the earlier photo-trained 3-head MLP (\vfour{}), whose top pose is the answer. The \emph{gate} is a newer 3-head MLP of the same architecture trained on the cue-verified video table plus training photos and photos of other poses (labelled transition/unknown); it answers ``is this one of my poses at all?'' and supplies the form score. The rule is fixed in advance: the answer is \vfour{}'s top pose unless the gate's top class is transition/unknown, in which case the system says it does not recognise a pose. This is a reject option in the sense of selective classification \cite{geifman2017selective}, applied at the system level. The policy was chosen on a validation half of the held-out photos and confirmed on the untouched half (Section~\ref{sec:results}).

\subsection{Rule engine and per-joint cues}
Each pose has a table of angle bands (e.g.\ for Downward Dog: knee $\in[110^\circ,180^\circ]$, shoulder $\in[95^\circ,180^\circ]$, hip $\in[20^\circ,140^\circ]$). The deviation of a joint is the distance in degrees outside its band (0 if inside), and the rule score is $1-\overline{d}/45^\circ$ clipped to $[0,1]$. Because the learned deviation head turned out to be weak (Section~\ref{sec:results}), the per-joint cue and the colour of each limb in the skeleton overlay come from these bands wherever a pose has them (17 of the 19 poses): coral if a joint is more than $15^\circ$ outside, amber above $7^\circ$, sage otherwise (Fig.~\ref{fig:occ}a). Two poses, Tree and Lunge, have no bands, and for them the cascade falls back to the learned head, so a correct Tree can show a coral leg (Fig.~\ref{fig:occ}d). We found this while verifying the calibration feature and report it as a limitation (Section~\ref{sec:limits}).

\subsection{ST-GCN ``flow check''}\label{sec:stgcn}
The sequence model treats a 60-frame window as a graph signal on the 33-joint MediaPipe skeleton. Coordinates are centred on the pelvis and divided by the hip width. Each of three ST-GCN blocks \cite{yan2018stgcn} applies a graph convolution with the symmetric-normalised skeleton adjacency $\mathbf{A}_{\mathrm{norm}}=\mathbf{D}^{-1/2}(\mathbf{A}+\mathbf{I})\mathbf{D}^{-1/2}$ (a $1{\times}1$ projection followed by aggregation over neighbours) and a 9-tap temporal convolution, with GELU and a residual connection; channels grow $3\to64\to128\to256$, features are averaged over joints and time, and a two-layer head outputs eight target poses plus transition/unknown (Fig.~\ref{fig:models}b). The app shows a flow-check answer only if its confidence reaches 0.70 (0.55 for child's pose); otherwise the UI shows ``static check'', and the per-frame cascade alone names the pose and scores form.

\runin{The input-rate contract.} The model is trained on 60 consecutive native-rate frames ($\approx$2.4\,s at 25\,fps). An early version of the web app fed its buffer one frame per finished server round trip (at most 2\,fps), so its ``60 frames'' spanned 30--120\,s and the model looked broken. The app now fills a time-based buffer from every camera frame and resamples it to 25\,fps; Section~\ref{sec:results} quantifies why this matters.
